% Split from 8-16yang.tex by cut-sheet v2/v3 — content below is the source scene, header adjusted
\chapter{Manifesto of the Merriweather Circle}
% \section*{Sub-Maths: Manifesto of Merriweather circle}

\lettrine{W}{illard} leaned back in his chair, his expression thoughtful but charged with quiet conviction. Across the table, Lee Yang watched him with the measured gaze of someone who respected Willard’s intellect but was wary of his penchant for challenging orthodoxy.

\begin{figure}[H]
    \centering
    \includegraphics[width=\textwidth]{Illustrations/vign-willardYang.png}
    \caption{Willard not vibing it with Yang}
    \label{fig:pirallingprimes}
\end{figure}

“Lee,” Willard began, his voice deliberate, “you’ve always struck me as someone who isn’t afraid to think against the grain. That’s why I feel you’ll appreciate this. What I’ve been doing with number theory—these ‘sub-mathematical probings,’ as I call them—is my attempt to explore mathematics not as an immutable monolith, but as a dynamic landscape of patterns waiting to be observed, modeled, and understood.”

Yang raised an eyebrow. “Sub-mathematical? That’s a curious term. Are you suggesting the opposite of meta-mathematics?”

“Precisely,” Willard replied. “Meta-mathematics, as the Austrians and their Vienna Circle positivists conceived it, saw mathematics as an abstract, axiomatic edifice—detached from experience, immune to the flux of time and matter. Carnap even had the audacity to claim that mathematics isn’t a language for expressing thoughts but merely the logical syntax of language. To them, mathematics was a sterile game, bound by rules but devoid of meaning.”

Yang nodded slowly, his gaze sharpening. “And Gödel proved them wrong.”

“Exactly,” Willard said, his tone rising slightly. “Gödel demonstrated that even within their precious axiomatic systems, there would always be truths that couldn’t be proven—truths that transcended their mechanical frameworks. And here’s the irony: Gödel, this master of abstraction, also discovered a static solution to Einstein’s field equations—a timeless, immutable picture of spacetime that defied the mutability of the physical universe.”

Yang interjected, “A picture that banished time itself. A \textit{Weltbild}—a worldview of reality as eternal, unchanging.”

“Right!” Willard exclaimed. “But where Gödel sought immutability, I see the opposite in mathematics. The objects of set theory may appear timeless, as Gödel claimed, but mathematics itself evolves—its truths, its tools, its relevance to science. This is where I diverge from the axiomatic purists. I don’t see mathematics as a temple to chip away at, stone by stone, with logical chisels. I see it as data—a sprawling, infinite dataset that demands exploration, visualization, and pattern recognition.”

Yang studied him for a moment, then asked, “But why call it sub-mathematics? Doesn’t that diminish its significance?”

Willard shook his head. “Not at all. It’s sub-mathematics because it grounds itself in the mutable, the empirical, the temporal. Where meta-mathematics floats above, detached and axiomatic, sub-mathematics delves below—into the chaotic, messy data of reality. It embraces the imperfections of computation, the contingency of models, the evolving nature of mathematical inquiry itself. It’s the complement to science, bound to time and mutability.”

Yang’s expression softened, as though he were beginning to see the contours of Willard’s vision. “You’re saying that sub-mathematics embraces the fleeting and the dynamic, while meta-mathematics chases the eternal and the unchanging.”

“Exactly,” Willard said, his voice almost reverent. “And isn’t that the ultimate irony of Gödel? In proving the limits of axiomatic systems, he paved the way for a mathematics that isn’t bound by those limits—a mathematics that can evolve, that can mirror the very flux of the universe. It’s not about rejecting immutability but complementing it with the mutable. Think of it as Planck meeting Gödel, quantum meeting continuum.”

Yang leaned back, his smile faint but genuine. “You make a compelling case, Willard. Perhaps there’s more to number theory than I gave it credit for. Maybe mathematics, like physics, needs its empiricists.”

% Willard chuckled softly. “And maybe physics, like mathematics, needs its dreamers.”

% \subsection*{Inter-Mutability in Physics and Mathematics}

Willard’s face lit up as he continued, “There’s something Roger Penrose once said about time that I think ties this all together. In physics, only objects with mass ‘experience’ time. And this experience is connected to the famous relationship between energy, mass, and time. Let me show you.”
He turned over his knapkin:
\[
E = hf= \frac{h}{T},
\]
where \(E\) is energy, \(h\) is Planck’s constant, and \(f\) is the frequency and using the relationship \(f = \frac{1}{T}\), where \(T\) is the period of oscillation. Now, equating two expressions for energy through Einstein’s mass-energy equivalence, \(E = mc^2\), and solving for \(T\) we find:
\[
mc^2 = \frac{h}{T}, \quad 
T = \frac{h}{mc^2}.
\]

“This equation,” Willard explained, “shows that the ‘experience’ of time—its duration—is tied directly to mass. Only objects with mass are subject to time’s mutability. Mathematics, in its immutable forms, doesn’t experience time. But what if we explore the mutable side of mathematics? That’s what I mean by sub-mathematics.”

Yang stared at the equation, nodding slowly. “So you’re saying that sub-mathematics brings mutability to an otherwise static discipline. A kind of temporal mathematics.”

“Exactly,” Willard said. “And in doing so, it gives mathematics a way to relate to the physical universe—not just as a tool, but as a dynamic partner in understanding reality.”

% \subsection*{Sub-Mathematics: A Grounding Manifesto}

Yang leaned back, the faint flicker of amusement crossing his face. "You know, Willard, if we're coining terms for this... perspective of yours, why not dabble with something a bit more... poetic? We've got \textit{meta-}, which means 'beyond' or 'about,' so maybe the antonym could be something like \textbf{intra-}, to signify 'within' mathematics."

Willard smiled faintly, shaking his head. "Intra- is tempting, but it feels... enclosed, doesn't it? This isn't about being constrained. It's about diving beneath the surface to what underlies mathematics. Something deeper."

Lee tilted his head. "Fair. Then what about \textbf{infra-}? It has that sense of being beneath, like in 'infrastructure.' A hidden, supporting layer."

Willard paused thoughtfully but shook his head again. "Close, but 'infra-' feels... functional. Like scaffolding or a wavelength band. We're not just dealing with hidden mechanics here; we're talking about a grounding philosophy—something that underscores mathematics itself."

Lee tapped his pen against the edge of his notebook. "How about \textbf{hypo-}? It has that same sense of 'beneath' but a more conceptual feel, like 'hypothesis' or 'hypothetical.'"

Willard gave a soft laugh. "Hypo- sounds too tentative, like we're theorizing. This isn't about speculation. It's about \textit{what mathematics stands on}, not some conjectural layer."

Finally, Willard leaned forward, his tone turning resolute. "\textbf{Sub-} is the word. It's not just beneath—it grounds mathematics. It’s the opposite of meta- in the clearest sense. Where meta-mathematics abstracts away from mathematics to examine its structure, \textbf{sub-mathematics} delves into its foundations. Think of mappings in mathematics—\textbf{injective} and \textbf{surjective}. Injection maps \textit{into} a structure, while surjection -as the French\footnote{French \textit{sur}=on of English.} would have us know- maps \textit{onto} it. Sub-mathematics is like an 'injection' into the grounding truths beneath mathematics itself."

Lee nodded slowly. "I can see that. \textbf{Sub-} has a clarity the others lack. It’s not about abstraction, it’s about anchoring. A kind of... foundation-laying for numbers."

% Willard smiled, pleased. "Exactly. Let’s not build castles in the air, Lee. Let’s anchor them in the bedrock of the earth."
